% Generated draft: requires mathematical and source review.
\documentclass[11pt]{article}
\usepackage[margin=25mm]{geometry}
\usepackage{fontspec}
\setmainfont{Latin Modern Roman}
\usepackage{amsmath,amssymb,mathtools,unicode-math}
\setmathfont{Latin Modern Math}
\usepackage{enumitem,needspace,xurl,hyperref}
\setlength{\parindent}{0pt}
\setlength{\parskip}{5pt}
\setlength{\emergencystretch}{3em}
\newcommand{\N}{\mathbb{N}}
\newcommand{\Z}{\mathbb{Z}}
\newcommand{\Q}{\mathbb{Q}}
\newcommand{\R}{\mathbb{R}}
\newcommand{\C}{\mathbb{C}}
\newcommand{\F}{\mathbb{F}}
\newcommand{\A}{\mathbb{A}}
\newcommand{\PP}{\mathbb{P}}
\newcommand{\E}{\mathbb{E}}
\newcommand{\eps}{\varepsilon}
\newcommand{\dd}{\,\mathrm{d}}
\begin{document}
\section*{UnsolvedMath ID: 30006440}
\textbf{Source problem code:} catalogue entry
\subsection{Definitions and mathematical statement}
Let \(n\geq 1\) be an integer, and let \(S_n\) be the group of permutations of \(\{1,\ldots,n\}\). For a subset \(A\subseteq S_n\) satisfying \(\langle A\rangle=S_n\), define the undirected Cayley graph \(\Gamma_n(A)\) to have vertex set \(S_n\), with distinct vertices \(g,h\) adjacent precisely when \(h=gs\) for some \(s\in A\cup A^{-1}\), where \(A^{-1}=\{a^{-1}:a\in A\}\). Thus inverses are available as single steps. Let \(d_A(g,h)\) be the minimum number of edges in a path from \(g\) to \(h\), and define
\[
\operatorname{diam}\Gamma_n(A)=\max_{g,h\in S_n}d_A(g,h),\qquad
D(n)=\max_{\substack{A\subseteq S_n\\\langle A\rangle=S_n}}\operatorname{diam}\Gamma_n(A).
\]
The singleton graph for \(n=1\) has diameter zero. Replacing \(A\) by \(A\cup A^{-1}\) does not change this graph, so this formulation also covers the convention of requiring an inverse-closed generating set.

\textbf{Proposed diameter bound.} The supplied statement asks to prove
\[
\forall n\geq 1,\qquad D(n)\leq \frac{1}{2}n^2+4n.
\]
Equivalently, the displayed upper bound is to hold for every generating subset of \(S_n\), with distance measured in the undirected graph defined above. The positive-integer range makes explicit the supplied phrase ``for every \(n\)''; the record gives no separate lower threshold.

\textbf{Proposed extremal families: incomplete specification.} The record also asks to prove or disprove that certain explicit generating families, described as related to involutions in a ``square-with-whiskers'' pattern, maximize the diameter for every \(n\). If \(W_n\subseteq S_n\) denotes one intended generating set, its proposed extremality would mean
\[
\langle W_n\rangle=S_n,
\qquad
\operatorname{diam}\Gamma_n(W_n)=D(n).
\]
This is only a schema: the supplied record contains no permutations defining \(W_n\), no construction of the pattern, and no parameter or admissible-size specifications for the plural families. It therefore does not determine a precise extremality conjecture or whether every member, or a selected member, is intended to attain the maximum. These data must be supplied before that part can be fully formulated.

The excerpt reports extensive searches, exhaustive in some cases, supporting the proposed extremality for \(n\leq 15\); it does not specify which cases were exhaustive. This is supplied computational evidence, not a proof for all \(n\). The excerpt also contains incomplete references to a question attributed to ``V.'' and a nilpotent-group conjecture involving ``J.''; neither supplies enough information to formulate an additional mathematical target.
\subsection{Short English statement}
The problem proposes that every undirected Cayley graph generated by permutations of \(n\) symbols has diameter at most \(\frac12 n^2+4n\). It also proposes that particular ``square-with-whiskers'' generating sets attain the largest possible diameter, but their construction is missing from the supplied record. Reported computations support extremality through \(n=15\), with exhaustiveness asserted only for some cases.
\subsection{Sources}
\begin{itemize}
\item Supplied local UnsolvedMath registry record, permanent ID \(30006440\), title ``Extremal Diameters of Symmetric-Group Cayley Graphs,'' exact source problem code \texttt{catalogue entry}, in the Oberwolfach Reports Open Problems collection. The record gives proposed year 2025 and labels the problem open. Its statement status is \texttt{corrected\_verified}, qualified by the supplied explanation that only formula delimiters, notation, grammar, and typography were repaired. These are registry labels, not independent mathematical verification.
\item Supplied report citation: \emph{MATRIX-MFO Tandem Workshop: Machine Learning and AI for Mathematics (2026)}, Oberwolfach Reports, DOI 10.4171/owr/2025/43, pp.\ 2323--2350. Supplied source URL and literature-source label ``Oberwolfach Report 43/2025'': \url{https://doi.org/10.4171/owr/2025/43}. The workshop year in the supplied title differs from the report identifier and proposed year; this discrepancy is unresolved.
\item Supplied literature source: ``CayleyPy Growth,'' \url{https://arxiv.org/abs/2509.19162}.
\item Attribution in the supplied excerpt: Alexander Chervov, in a contribution titled ``CayleyPy -- Python AI-based library for google size Cayley graphs,'' described as joint work with A. Soibelman, S. Galkin, D. Fedoriaka, E. Konstantinova, A. Naumov, I. Kiselev, A. Sheveleva, I. Koltsov, S. Lytkin, A. Smolensky, F. Levkovich-Maslyuk, R. Grimov, D. Volovich, H. Isambert, et al. The registry's separate proposer field is empty; the excerpt does not establish individual authorship of either conjecture.
\item Literature assessment dated exactly \texttt{2026-08-22}: the supplied curation labels the refined bound and square-with-whiskers extremality open and explicitly conjectural, and describes computational evidence confined to the reported finite range \(n\leq 15\). This is supplied dated curation, not an independently verified literature assessment. Neither linked source was independently consulted in preparing this draft.
\end{itemize}
\end{document}
